% !PRINTED 106
\weberpage{106}
%<<CONT

\[
 \varphi(h)=h f'(x)+\frac{h^2}{1\cdot2}f''(x)+\cdots+h^n a_0
\]
\[
 |f(x)|-|\varphi(h)|<|f(x+h)|<|f(x)|+|\varphi(h)|,
\]
or
\[
 -|\varphi(h)|<|f(x+h)|-|f(x)|<|\varphi(h)|.
\]
Now, if $\varrho<\varepsilon$, then $|\varphi(h)|<\omega$, and we can therefore also say:

\medskip
\noindent VI. \gesperrt{The absolute value of the difference}
\[
 |f(x+h)|-|f(x)|
\]
\gesperrt{is smaller than an arbitrarily given magnitude $\omega$ as soon as the absolute value of $h$ is smaller than a sufficiently small $\varepsilon$.}
\medskip

From this we can still conclude, if we put $x=y+iz$ and assume $h$ either real or purely imaginary, that $|f(y+zi)|$ is, for fixed $z$, a continuous function of $y$, and for fixed $y$, a continuous function of $z$.

\medskip
\noindent VII. \gesperrt{If the coefficients $a_0,a_1,\ldots,a_n$ and the variable $x$ are real, then $x$ may be chosen, in absolute value, so large that the sign of $f(x)$ agrees with the sign of the first term $a_0x^n$; thus, for positive $a_0$ and even $n$, it becomes positive, while for odd $n$ it becomes negative for negative $x$ and positive for positive $x$.}

For in
\[
 f(x)=x^n\left(a_0+\frac{a_1}{x}+\frac{a_2}{x^2}+\cdots+\frac{a_n}{x^n}\right)
\]
one can choose $x$ so large in absolute value that the absolute value of
\[
 \frac{a_1}{x}+\frac{a_2}{x^2}+\cdots+\frac{a_n}{x^n}
\]
falls below that of $a_0$, and therefore the expression in parentheses
\[
 a_0+\frac{a_1}{x}+\frac{a_2}{x^2}+\cdots+\frac{a_n}{x^n}
\]
has the sign of $a_0$, which it then retains for every absolutely larger $x$.
